\documentclass[12pt]{article}

% Idioma e codificação
\usepackage[brazil]{babel}
\usepackage[utf8]{inputenc}
\usepackage[T1]{fontenc}

% Matemática e listas
\usepackage{amsmath, amssymb}
\usepackage{enumitem}

% Layout e cores
\usepackage{geometry}
\usepackage{graphicx}
\usepackage{xcolor}

% Cores institucionais
\definecolor{maincolor}{RGB}{0,100,50}
\definecolor{forestgreen}{RGB}{0,100,50}

% Cabeçalho e rodapé
\usepackage{fancyhdr}

% Blocos
\usepackage[most]{tcolorbox}
\newtcolorbox{block}[1]{%
  colback=green!5,
  colframe=maincolor,
  coltitle=white,
  title=#1,
  fonttitle=\bfseries,
  boxrule=1.5pt,
  arc=3mm,
  left=6pt,
  right=6pt,
  top=6pt,
  bottom=6pt
}

% Margens
\geometry{margin=2.5cm}

% Fancyhdr
\pagestyle{fancy}
\fancyhf{}

\renewcommand{\headrulewidth}{2pt}
\renewcommand{\footrulewidth}{2pt}

\renewcommand{\headrule}{\hbox to\headwidth{%
  \color{forestgreen}\leaders\hrule height \headrulewidth\hfill}}

\renewcommand{\footrule}{\hbox to\headwidth{%
  \color{forestgreen}\leaders\hrule height \footrulewidth\hfill}}

% Altura do cabeçalho (corrigida)
\setlength{\headheight}{52pt}
\addtolength{\topmargin}{-20pt}

% Cabeçalho
\fancyhead[L]{%
  \includegraphics[height=1.2cm]{Universidade_de_Rio_Verde.jpg}
}

\fancyhead[R]{%
  \raggedleft
  \textcolor{forestgreen}{%
    \small
    \textbf{Lista de Exercícios}\\
    Me.~Sergio Souza Novak\\
    \textit{Processo de Software}
    }
    }
    
    % Rodapé
    \fancyfoot[C]{\textcolor{forestgreen}{\thepage}}
    
    % =========================
    \begin{document}
    % =========================
    
    \vspace*{0.2cm}
    
\begin{center}
  {\Large \textbf{Lista de Exercícios: Organização de Pastas e Arquivos}}\\[0.2cm]
  {\normalsize Informática Básica --- Criação de Estruturas de Diretórios}
\end{center}

\vspace{0.4cm}

\textbf{Objetivo:} Praticar a criação de pastas, subpastas e a organização de arquivos com diferentes extensões no sistema operacional.

\textbf{Instruções Gerais:}
\begin{itemize}
    \item Crie no seu computador uma pasta principal com o seu \textbf{Nome Completo} na Área de Trabalho.
    \item Dentro dela, monte as estruturas de pastas e arquivos conforme indicado em cada exercício.
    \item Para criar os arquivos, você pode criar documentos em branco (txt, docx ou xlsx) salvando-os com os nomes e extensões indicadas.
    \item Responda às questões propostas.
\end{itemize}

\vspace{0.3cm}

% -------------------------------------------------------
\begin{block}{Exercício 1 --- Empresa Sementeira}
% -------------------------------------------------------

A empresa \textbf{SementeVerde Agro} precisa organizar seus arquivos de gestão e controle de estoque. Crie no seu computador a seguinte estrutura:

\begin{verbatim}
SementeVerde Agro
|-- Administrativo
|   |-- contrato_arrendamento.docx
|   \-- cadastro_funcionarios.xlsx
|-- Comercial
|   |-- lista_clientes.xlsx
|   \-- tabela_precos_2026.docx
\-- Estoque
    |-- controle_soja.xlsx
    \-- laudo_milho.docx
\end{verbatim}

\textbf{Questões:}
\begin{enumerate}[label=\alph*)]
    \item Quantas pastas e quantos arquivos você criou neste exercício?
    \item Em qual pasta está salvo o arquivo \texttt{tabela\_precos\_2026.docx}?
    \item Quais são as extensões dos arquivos utilizados no exercício?
\end{enumerate}

\end{block}

\vspace{0.4cm}

% -------------------------------------------------------
\begin{block}{Exercício 2 --- Empresa Madeireira}
% -------------------------------------------------------

A empresa \textbf{MadeiraForte SA} precisa organizar seus documentos de manejo florestal, produção e vendas. Monte a seguinte estrutura:

\begin{verbatim}
MadeiraForte SA
|-- Manejo Florestal
|   \-- Licencas
|       \-- IBAMA
|           \-- licenca_ambiental.docx
|-- Producao
|   \-- Corte
|       \-- Talhao 01
|           |-- plano_de_corte.xlsx
|           \-- mapa_plantio.png
\-- Vendas
    \-- Atacado
        \-- Nacional
            |-- orcamento_moveis.docx
            \-- pedido_aprovado.docx
\end{verbatim}

\textbf{Questões:}
\begin{enumerate}[label=\alph*)]
    \item Qual é o caminho completo de pastas para encontrar o arquivo \texttt{licenca\_ambiental.docx}?
    \item Quantos níveis de subpastas existem na seção de \textbf{Vendas} até chegar ao arquivo \texttt{orcamento\_moveis.docx}?
    \item Se você adicionar um arquivo chamado \texttt{nota\_fiscal.docx} dentro da pasta \textbf{Nacional}, em qual nível de profundidade ele estará localizado?
\end{enumerate}

\end{block}

\vspace{0.4cm}

% -------------------------------------------------------
\begin{block}{Exercício 3 --- Empresa de Consórcio Agrícola}
% -------------------------------------------------------

A empresa \textbf{Consórcio Campo Unido} gerencia dados de associados e contratos de consórcio. Crie a estrutura de diretórios e arquivos a seguir:

\begin{verbatim}
Consorcio Campo Unido
|-- Associados
|   |-- cadastro_produtores.xlsx
|   \-- documentos_joao.docx
|-- Consorcio Tratores
|   |-- contrato_grupo_A.docx
|   \-- recibo_parcela.docx
\-- Financeiro
    |-- arrecadacao_mensal.xlsx
    \-- prestacao_contas.docx
\end{verbatim}

\textbf{Questões:}
\begin{enumerate}[label=\alph*)]
    \item Onde está armazenado o contrato do grupo de tratores?
    \item Qual programa você usaria para abrir o arquivo \texttt{arrecadacao\_mensal.xlsx}?
    \item Se fosse adicionado um consórcio de colheitadeiras, qual pasta você criaria e qual arquivo colocaria dentro dela?
\end{enumerate}

\end{block}

% =========================
\end{document}
% =========================
